\section{Numerically Verified Piecewise Drift Candidate}
\label{sec:appendix_piecewise}

The negative results in Table~\ref{tab:sruas_audit} demonstrate that the simple SRUAS sufficient condition is too conservative for the benchmark default, but they do not rule out the existence of alternative Lyapunov constructions.
The checks reported in this appendix provide \textbf{computationally verified negative drift on exhaustive boundary shells} at $(\calbeta,\calgamma,\caloffset)=(0.85,0.5,0.5)$, constituting strong evidence for a viable piecewise Lyapunov proof route, though they do not constitute a formal symbolic proof.
To probe the gap, we performed a separate post-validation search over convex combinations of polynomial, imbalance, and smoothed-max Lyapunov templates, with an exact-shell refinement loop that feeds boundary counterexamples back into the search.

For the benchmark-default reflected policy $(\calbeta,\calgamma,\caloffset)=(0.85,0.5,0.5)$ at $\temp=20$, the best refined candidate is
\begin{equation}
\label{eq:piecewise_candidate}
\begin{aligned}
\lyap_{\mathrm{PW}}(\Qstate)
&=
0.541679
\left[
\frac12 \sum_{i=1}^{\Nservers} \frac{\Qlen{i}^2}{\srvrate_i}
+
2 \sum_{i=1}^{\Nservers}
\left(
\frac{\Qlen{i}}{\srvrate_i}
-
\frac{1}{\Nservers}\sum_{j=1}^{\Nservers}\frac{\Qlen{j}}{\srvrate_j}
\right)^2
\right] \\
&\quad
+ 0.187796 \left(\sum_{i=1}^{\Nservers}\Qlen{i}\right)^2
+ 0.175661 \sum_{i=1}^{\Nservers}\frac{\Qlen{i}^3}{\srvrate_i}
+ 0.092741 \left(\sum_{i=1}^{\Nservers}\frac{\Qlen{i}}{\srvrate_i^{1/2}}\right)^2
+ 0.002124 \max_{1\le i\le \Nservers}\Qlen{i}^2.
\end{aligned}
\end{equation}
where the displayed coefficients are rounded to six decimal places from the verification capsule.

These coefficients are benchmark-specific; they serve as existence evidence only. \cref{eq:piecewise_candidate} is a \emph{numerically selected} benchmark-specific candidate, not a theorem-backed symbolic Lyapunov function.
Its role is evidentiary: it shows that the failure of $\varepsilon_{\calbeta,\calgamma,\caloffset}>0$ does not prevent the existence of much stronger negative-drift candidates near the benchmark tail boundary.

The dedicated verification capsule evaluates the exact generator of the benchmark-default reflected policy against \cref{eq:piecewise_candidate} on three exhaustive boundary shells and on additional high-load random shells. Note that states inside the candidate compact set $C$ (where $|\Qstate|_1 \le 15$) are structurally excluded from verification, as strict negative drift is only required outside $C$, while internal drift is accommodated by the uniform bound $R$.
Table~\ref{tab:piecewise_candidate} reports the results.

\begin{table}[htbp]
\centering
\small
\caption{Numerical piecewise-drift verification for the benchmark-default reflected policy at $\temp=20$. The candidate \eqref{eq:piecewise_candidate} is \emph{not} a formal proof, but its exact generator evaluation is strictly negative on all states in the three verified boundary shells and on all sampled higher-load shells listed below.}
\label{tab:piecewise_candidate}
\renewcommand{\arraystretch}{1.2}
\begin{tabular}{@{}l r r r c@{}}
\toprule
\textbf{Check set} & \textbf{States} & \textbf{Max drift} & \textbf{Positive states} & \textbf{Pass?} \\
\midrule
$|\Qstate|_1 = 16$ & 2{,}042{,}975 & $-4.6446$ & 0 & Yes \\
$|\Qstate|_1 = 17$ & 3{,}124{,}550 & $-4.6447$ & 0 & Yes \\
$|\Qstate|_1 = 18$ & 4{,}686{,}825 & $-4.9245$ & 0 & Yes \\
Random shells $|\Qstate|_1 \in \{24,32,48,64,96,128,160,256\}$ & 8{,}192 & $-21.5534$ & 0 & Yes \\
\bottomrule
\end{tabular}
\end{table}

These checks substantially strengthen the empirical-theory interface for the benchmark-default reflected point.
However, this does not mathematically prove a full positive Harris recurrence theorem for the unrestricted reflected family because it does not provide a symbolic drift inequality valid for \emph{all} sufficiently large states.
Likewise, the verified shells do not constitute an exhaustive finite-state audit of an entire compact region, so they cannot by themselves be combined with the current analytic bounds into a formal Foster--Lyapunov certificate.
Accordingly, the manuscript treats \cref{eq:piecewise_candidate} and Table~\ref{tab:piecewise_candidate} purely as numerical evidence for a promising piecewise Foster--Lyapunov route, not as a formal replacement for the certified guarantees of \cref{thm:sruas}.
